\chapter{Applications of Control Theory to Biological Systems}
\label{ch:control-bio}

\epigraph{%
  The body is not a robot that happens to be made of meat.
  It is a system that has evolved to exploit its own physics.
}{}

\section{Forward Dynamics Simulation}

Forward dynamics simulation integrates the equations of motion given muscle
excitations $\bm{u}(t)$:
\begin{equation}
  \ddot{\bm{q}} = M(\bm{q})^{-1}\left[R(\bm{q})\bm{F}_{\text{muscle}} + J_c^T \bm{F}_c
  - C(\bm{q}, \dot{\bm{q}})\dot{\bm{q}} - g(\bm{q})\right],
  \label{eq:ch10:forward-dynamics}
\end{equation}
coupled with the muscle state equations (Chapter~3):
\begin{align}
  \dot{a}_i &= (u_i - a_i) / \tau_i, \\
  \dot{\ell}_i^M &= v_i^M(a_i, \ell_i^M, \ell_i^{MT}(\bm{q})).
\end{align}

\begin{lstlisting}[language=Python, caption={Forward dynamics integration for a musculoskeletal model.}]
import numpy as np
from scipy.integrate import solve_ivp
from typing import Callable

def forward_dynamics_simulation(
    mass_matrix_func: Callable,
    coriolis_func: Callable,
    gravity_func: Callable,
    moment_arm_func: Callable,
    muscle_force_func: Callable,
    excitation_func: Callable,
    q0: np.ndarray,
    qdot0: np.ndarray,
    a0: np.ndarray,
    t_span: tuple[float, float],
    n_eval: int = 500
) -> dict:
    """Simulate forward dynamics of a musculoskeletal system.

    Parameters
    ----------
    mass_matrix_func : callable
        q -> M(q), shape (n_dof, n_dof)
    coriolis_func : callable
        (q, qdot) -> C(q, qdot) @ qdot, shape (n_dof,)
    gravity_func : callable
        q -> g(q), shape (n_dof,)
    moment_arm_func : callable
        q -> R(q), shape (n_dof, n_muscles)
    muscle_force_func : callable
        (a, l_M, v_M) -> F_muscle for each muscle
    excitation_func : callable
        t -> u(t), shape (n_muscles,)
    q0 : np.ndarray
        Initial joint angles.
    qdot0 : np.ndarray
        Initial joint velocities.
    a0 : np.ndarray
        Initial muscle activations.
    t_span : tuple
        (t_start, t_end)
    n_eval : int
        Number of output time points.

    Returns
    -------
    dict
        Keys: 't', 'q', 'qdot', 'a', 'tau'
    """
    n_dof = len(q0)
    n_muscles = len(a0)
    t_eval = np.linspace(t_span[0], t_span[1], n_eval)

    def dynamics(t: float, state: np.ndarray) -> np.ndarray:
        q = state[:n_dof]
        qdot = state[n_dof:2*n_dof]
        a = np.clip(state[2*n_dof:], 0.0, 1.0)

        # Muscle forces (simplified: assume known fiber length)
        F_muscle = np.array([
            a[i] * 500.0  # Simplified: F = a * F0
            for i in range(n_muscles)
        ])

        # Joint torques
        R = moment_arm_func(q)
        tau = R @ F_muscle

        # Skeletal dynamics
        M = mass_matrix_func(q)
        C_qdot = coriolis_func(q, qdot)
        g = gravity_func(q)
        qddot = np.linalg.solve(M, tau - C_qdot - g)

        # Activation dynamics
        u = excitation_func(t)
        tau_act = 0.015  # 15 ms activation
        tau_deact = 0.050  # 50 ms deactivation
        adot = np.where(
            u > a,
            (u - a) / tau_act,
            (u - a) / tau_deact
        )

        return np.concatenate([qdot, qddot, adot])

    state0 = np.concatenate([q0, qdot0, a0])
    sol = solve_ivp(dynamics, t_span, state0, t_eval=t_eval,
                    method='RK45', max_step=0.001)

    return {
        't': sol.t,
        'q': sol.y[:n_dof].T,
        'qdot': sol.y[n_dof:2*n_dof].T,
        'a': sol.y[2*n_dof:].T,
    }
\end{lstlisting}

\section{Predictive Simulation}

Predictive simulation uses trajectory optimization (Volume~I, Chapter~5) to find
the muscle excitations that produce a desired motion while minimizing a cost:
\begin{equation}
  \min_{\bm{u}(\cdot)} \int_0^T \mathcal{L}(\state, \bm{u}) \dt
  \quad \text{s.t.} \quad \dot{\state} = f(\state, \bm{u}),\;
  \state(0) = \state_0,\; h(\state, \bm{u}) \leq 0,
  \label{eq:ch10:predictive-sim}
\end{equation}
where $\mathcal{L}$ is a physiologically motivated cost (e.g., metabolic energy,
muscle fatigue, deviation from desired kinematics) and $h$ captures physical
constraints (joint limits, muscle force capacity, contact constraints).

\noindent
Agrachev \& Sachkov \cite{AgrachevSachkov2004}, Theorem~8.1 (Pontryagin's Maximum Principle),
establishes the fundamental existence and optimality conditions for such optimal control problems.
Their treatment covers affine systems with constraints and provides both theoretical guarantees
and computational methods (the method of characteristics) for solving OCPs of this form.

\subsection{Metabolic Cost Models}

Predictive simulation requires a physiologically meaningful running cost. The
most widely used choice in gait simulation is a \emph{metabolic energy rate} model,
which converts muscle state and force into an estimate of the whole-body
rate of ATP consumption. Umberger et al.\ \cite{umberger2003metabolic} proposed a
four-term decomposition in which the instantaneous metabolic rate of muscle $i$
is written as
\[
  \dot{E}_i = \dot{h}_i^{A} + \dot{h}_i^{M} + \dot{h}_i^{SL} + \dot{w}_i^{CE},
\]
summing activation heat, maintenance heat, shortening/lengthening heat, and the
mechanical work done by the contractile element, each calibrated against in-vitro
muscle energetics. Bhargava et al.\ \cite{bhargava2004metabolic} offer a closely
related formulation that separates activation and maintenance using fast and slow
fiber fractions, and is the model shipped with OpenSim
\cite{delp2007opensim,rajagopal2016full}. Whole-body metabolic cost is obtained by
summing $\dot{E}_i$ across all muscles and adding a constant basal rate.

In practice, minimizing $\int_0^T \sum_i \dot{E}_i \dt$ is expensive and nonsmooth,
so many optimal-control studies substitute a surrogate such as $\sum_i a_i^p$ or
$\sum_i (F_i/F_i^0)^p$ with $p \in \{2, 3\}$
\cite{pandy2001computer,anderson2001dynamic}. Surrogates and calibrated metabolic
models generally agree on qualitative patterns (preferred walking speed, stride
frequency) but can disagree on absolute cost.

\subsection{Muscle Synergies}

Predictive simulation of a 600-muscle model has many more inputs than degrees of
freedom, reproducing the motor-redundancy problem of Chapter~1 at the control
level. A long line of neuroscience research has proposed that the central nervous
system sidesteps this curse of dimensionality by driving muscles in low-dimensional
patterns called \emph{muscle synergies}: a small set of time-invariant activation
vectors $\bm{w}_k \in \Reals^{n_{\text{m}}}$ modulated by time-varying coefficients
$c_k(t)$ so that the full activation is
\[
  \bm{a}(t) \approx \sum_{k=1}^{K} c_k(t)\, \bm{w}_k,
  \qquad K \ll n_{\text{m}}.
\]
Bizzi \& Cheung \cite{bizzi2008combining} summarize the evidence for synergies as
building blocks of motor behavior across species, while d'Avella et al.\
\cite{davella2003combinations} showed that a handful of time-varying synergies
account for most of the EMG variance in natural arm reaching. Non-negative matrix
factorization \cite{HarrisWolpert1998} is the standard extraction technique because
activations and synergies are non-negative.

For control-theoretic purposes, synergies provide a rationale for low-dimensional
parameterizations of the input in optimal-control problems, reducing the effective
control dimension from $n_{\text{m}}$ to $K$ and making predictive simulation
tractable at whole-body scale.

\section{Prosthetics and Exoskeleton Control}

The control theory of Volumes~I--II applies directly to assistive devices:
\begin{itemize}
  \item \textbf{Powered prostheses}: trajectory tracking with impedance modulation
  \item \textbf{Exoskeletons}: assist-as-needed control using contraction theory
  \item \textbf{FES (Functional Electrical Stimulation)}: direct muscle activation
        control for paralyzed limbs
\end{itemize}

\section{Sport Biomechanics Applications}

\subsection{The Golf Swing as a Biomechanical System}

The golf swing demonstrates every concept in this volume:
\begin{enumerate}
  \item Multi-segment kinematic chain (15+ DOF)
  \item Muscle-driven actuation with co-contraction
  \item Bi-articular coupling (forearm muscles, lats)
  \item Passive dynamics exploitation (wrist release)
  \item Shaft flexibility (deformable body)
  \item Ground reaction force coupling (closed chain at the feet)
\end{enumerate}

The complete biomechanical analysis integrates:
\begin{itemize}
  \item Motion capture (kinematics)
  \item Force plates (ground reaction forces)
  \item EMG (muscle activation patterns)
  \item Inverse dynamics (joint torques)
  \item Static optimization (muscle forces)
  \item Forward dynamics validation (predictive simulation)
\end{itemize}

\section{Synthesis: From Theory to Application}

The progression through this volume mirrors the progression through the series:

\begin{center}
\renewcommand{\arraystretch}{1.3}
\begin{tabular}{@{}llll@{}}
\toprule
\textbf{Volume} & \textbf{Concept} & \textbf{Biomechanical Application} & \textbf{Agrachev \& Sachkov} \\
\midrule
Vol~0   & Screw axes, PoE      & Joint kinematics, ISA analysis            & \cite[Ch.~2]{AgrachevSachkov2004} \\
Vol~I   & Contraction theory   & Stability of muscle-driven systems        & \cite[Ch.~4]{AgrachevSachkov2004} \\
Vol~I   & DDP/iLQR             & Optimal muscle excitation patterns        & \cite[Ch.~8]{AgrachevSachkov2004} \\
Vol~II  & Trajectory tubes     & Motor variability envelopes               & \cite[Ch.~3]{AgrachevSachkov2004} \\
Vol~II  & Phase variables      & Gait cycle parameterization               & \cite[Ch.~7]{AgrachevSachkov2004} \\
Vol~III & Hill model           & Force prediction                          & \cite[Ch.~6]{AgrachevSachkov2004} \\
Vol~III & Inverse dynamics     & Joint torque estimation                   & \cite[Ch.~5]{AgrachevSachkov2004} \\
Vol~IV  & Motor control        & Neural command generation                 & \cite[Ch.~8]{AgrachevSachkov2004} \\
\bottomrule
\end{tabular}
\end{center}

\section*{Exercises}

\begin{enumerate}
  \item \textbf{Forward Simulation.}
        Using the provided code (adapted for a simple 2-DOF model), simulate a
        reaching movement with prescribed muscle excitation ramps.

  \item \textbf{Predictive vs.\ Tracking.}
        Compare the muscle activations from static optimization (tracking measured
        kinematics) with those from predictive simulation (optimizing a metabolic cost).

  \item \textbf{Prosthetic Control.}
        Design an impedance controller for a prosthetic knee that provides different
        stiffness during stance and swing phases.

  \item \textbf{(Programming.)} Implement a complete inverse dynamics $\to$ static
        optimization pipeline for a 3-DOF planar arm model with 6 muscles.
\end{enumerate}
